\documentclass[11pt,a4paper]{coverletter}

\senderblock
  {Radheem Bin Razi}
  {Distributed Systems \& Cloud-Native Engineer}
  {Ilmenau, Germany \textbar\ \href{mailto:sheikh.radheem@gmail.com}{sheikh.radheem@gmail.com} \textbar\ \href{https://radheem.github.io/my-cv}{portfolio}}

\begin{document}

%% ===== ENGLISH =====
\recipient{Darktrace}{Hiring Team}{}

\opening{Dear Hiring Team,}

A reliable data backbone quietly enables AI to spot threats before they cause damage, and designing that foundation is exactly why I want to join Darktrace as a Data Engineer. Your platform turns complex telemetry into proactive defense, and I want to help build the infrastructure that keeps those models accurate. Over five years, I have focused on distributed systems and cloud-native pipelines that scale without breaking under pressure.

I have spent recent years building event-driven ingestion systems that survive worker failures and keep data moving to AI models. I designed a durable pipeline that automatically retries each extraction and embedding step, streaming real-time progress through NATS JetStream while keeping all inference local. I later extended that reliability mindset to machine learning workflows by deploying a full training and serving lifecycle on Kubernetes. A configuration-driven Python client automated everything from feature group creation to model deployment, while a Kafka metrics pipeline fanned out telemetry to multiple databases for monitoring. These experiences gave me direct practice in keeping data flows predictable, which aligns closely with your need for scalable pipelines.

Cloud-native deployment and cross-team collaboration are equally important when data infrastructure powers production AI. I have packaged complex research frameworks into reproducible Kubernetes clusters using Helm, and I routinely document architectural decisions so data scientists and MLOps engineers can integrate smoothly. Working on a top-graded academic project, I wired together Cassandra, InfluxDB, and object storage to feed a prediction model, and I built a lightweight SDK that let non-infrastructure teammates trigger training jobs without touching the command line. This blend of infrastructure rigor and developer-friendly tooling helps me bridge the gap between raw data collection and model production.

I am ready to bring this operational discipline to your R\&D team in The Hague and would welcome the opportunity to discuss how my background aligns with your data platform goals. I am available to start full-time immediately, though I remain open to working-student or part-time arrangements. I am currently based in Germany and can relocate to the Netherlands within two to three weeks, and I am fully Blue Card-ready so your team only needs to issue the contract. Thank you for your time.

\closing{Sincerely,}{Radheem Bin Razi}

\clearpage
\selectlanguage{ngerman}

%% ===== DEUTSCH =====
\recipient{Darktrace}{Personalabteilung}{}

\opening{Sehr geehrtes Hiring-Team,}

Eine zuverlässige Datenbasis ermöglicht es der KI im Hintergrund, Bedrohungen zu erkennen, bevor sie Schaden anrichten. Genau das Entwerfen dieser Grundlage ist der Grund, warum ich als Data Engineer zu Darktrace wechseln möchte. Ihre Plattform wandelt komplexe Telemetriedaten in proaktive Abwehrmaßnahmen um, und ich möchte die Infrastruktur aufbauen, die die Genauigkeit dieser Modelle sicherstellt. In den vergangenen fünf Jahren habe ich mich auf verteilte Systeme und cloud-native Pipelines konzentriert, die auch unter hoher Last skalieren, ohne an Stabilität zu verlieren.

In den letzten Jahren habe ich ereignisgesteuerte Ingestion-Systeme aufgebaut, die Worker-Ausfälle überstehen und den Datenfluss zu den KI-Modellen aufrechterhalten. Ich habe eine robuste Pipeline entworfen, die jeden Extraktions- und Embedding-Schritt automatisch wiederholt, während der Echtzeitfortschritt über NATS JetStream gestreamt wird und alle Inference-Schritte lokal bleiben. Später habe ich dieses Prinzip der Zuverlässigkeit auf Machine-Learning-Workflows übertragen, indem ich einen vollständigen Trainings- und Serving-Lifecycle auf Kubernetes bereitgestellt habe. Ein konfigurationsbasierter Python-Client automatisierte dabei alles von der Erstellung von Feature-Gruppen bis zur Modellbereitstellung, während eine Kafka-Metriken-Pipeline Telemetriedaten zur Überwachung auf mehrere Datenbanken verteilte. Diese Erfahrungen haben mir direkte Praxis darin vermittelt, Datenflüsse stabil und vorhersehbar zu halten, was eng mit Ihrem Bedarf an skalierbaren Pipelines übereinstimmt.

Cloud-native Deployment und teamübergreifende Zusammenarbeit sind ebenso wichtig, wenn Dateninfrastruktur produktive KI-Systeme antreibt. Ich habe komplexe Forschungsframeworks mit Helm in reproduzierbare Kubernetes-Cluster überführt und dokumentiere architektonische Entscheidungen regelmäßig, damit Data Scientists und MLOps-Ingenieure nahtlos anbinden können. Im Rahmen eines sehr gut bewerteten akademischen Projekts habe ich Cassandra, InfluxDB und Object Storage verbunden, um ein Vorhersagemodell zu speisen, und ein leichtgewichtiges SDK entwickelt, das Teamkollegen ohne Infrastrukturhintergrund das Auslösen von Training-Jobs ohne Befehlszeile ermöglichte. Diese Kombination aus infrastrukturtechnischer Präzision und benutzerfreundlicher Tooling hilft mir, die Lücke zwischen roher Datenerfassung und Modellproduktion zu schließen.

Ich bin bereit, diese operative Disziplin in Ihr F\&E-Team in Den Haag einzubringen und begrüße die Gelegenheit, in einem Gespräch zu besprechen, wie mein Profil zu Ihren Zielen für die Datenplattform passt. Ich stehe ab sofort vollzeitlich zur Verfügung, bin jedoch auch offen für eine Tätigkeit als Werkstudent oder in Teilzeit. Ich befinde mich derzeit in Deutschland und kann innerhalb von zwei bis drei Wochen in die Niederlande umziehen. Ich bin vollständig für die Blue Card vorbereitet, sodass Ihr Team lediglich den Arbeitsvertrag ausstellen muss. Vielen Dank für Ihre Zeit.

\closing{Mit freundlichen Grüßen,}{Radheem Bin Razi}

\end{document}
